\section{Geometric Derivation of the Propagation Parameters}
\label{sec:strategy2-parameter-bridges}

This appendix derives the real parameters used in the finite Strategy~2 inequalities.  It is the only place where geometric data and the
self-contained optimization variables coexist.  Every proposition below has
one task: produce a point of an explicitly stated real domain and identify
the objective with the scalar quantity used in the body.  No optimization
inequality is proved here.

\subsection{Signed center and the four all-Vd0 problems}

\begin{proposition}[Signed-domain bridge]
\label{prop:s2-domain-bridge}
Let \((R,\alpha,\delta)\) be the normalized signed-center parameter tuple.  Under
\[
 r=R,\qquad a=\alpha,\qquad d=\delta,
\]
the full geometric CE1 and CE2 center cells are respectively
\(\mathcal D_{\rm C1}\) and \(\mathcal D_{\rm C2}\).  The additional
orientation and handoff inequalities in the CE1 all-Vd0 one-gap branch give
the smaller cell \(\mathcal D_1\), while
\(\mathcal D_2=\mathcal D_{\rm C2}\).
\end{proposition}

\begin{proof}
The signed-center normal form gives
\[
 W=1-R,\quad E=\sqrt{1-RW},\quad
 \eta=1-E,\quad P=E(1-E),
\]
\[
 k=\eta+\alpha+\delta,\quad
 \Delta_R=P-\alpha-W\delta,\quad
 \Delta_L=P-R\alpha-\delta.
\]
After substitution these are the definitions of
\(w,e,\eta,p,k,\Delta_R,\Delta_L\).  The CE sign is exactly the sign of
\(\Delta_L\).  In the oriented CE1 one-gap branch the midpoint and endpoint
orders additionally give
\[
 \alpha<W/2,\qquad \delta<R/2,\qquad
 R\delta-W\alpha>0,
\]
which are precisely the extra inequalities defining \(\mathcal D_1\).
\end{proof}

Every radial lower bound in the one-gap, paired-gap, and propagated
problems is strictly larger than $1/2$.  On this high-radial domain the exact
function dictionary is therefore
\[
 F(c,x)=g_c(1-x),
 \qquad
 T(c,x)=g_c^\vee(x).
\]
For the endpoint parameterss below,
\[
 \mathrm{Cap}(\mathbf e)
 =g_{c_-}(1-a_-)+g_{c_+}(1-a_+).
\]

\begin{proposition}[Endpoint bridge]
\label{prop:s2-endpoint-bridge}
The normalized one-gap endpoint parameters is
\[
 \left(\frac{k}{r},\ r-d,\ 1-\frac a w,\ 1-\frac d r\right),
\]
and
\[
 \mathrm{Cap}(\mathbf e_{\rm one})-1
 =\Psi_{\rm E1}(r,a,d).
\]
On \(\mathcal D_2\), the paired endpoint data are
\[
 \left(w-a,\ r-d,\ \frac{w-a}{w},\ \frac{r-d}{r}\right),
\]
and
\[
 \mathrm{Cap}(\mathbf e_{\rm two})-1
 =\Psi_{\rm E2}(r,a,d).
\]
\end{proposition}

\begin{proof}
Substitute the signed trace endpoints and complementary radial exits, then use
\(F(c,x)=g_c(1-x)\).
\end{proof}

\begin{proposition}[Returned-reach lower bound bridge]
\label{prop:s2-return-bridge}
Put
\[
 X=r-d,\qquad H=\frac{k}{2r},\qquad
 m=\min\left\{\frac ar,\frac dw\right\}.
\]
Let \(Y_0=H\) and
\[
 Y_1=T(1-a,Y_0),\qquad
 Y_2=T(1-m,Y_1),\qquad
 Y_3=T(1-d,Y_2).
\]
The objective \(Y_3>w+d=1-X\) in Problem S2-R1 or S2-R2 is the
complement-dual form of the middle-three exact forward-transfer inequality.
It implies the body target \(Z>1-H\).
\end{proposition}

\begin{proof}
The five radial lower bounds are
\[
 1-\frac dr,\quad 1-d,\quad 1-m,\quad 1-a,\quad 1-\frac aw.
\]
The first and fifth transfers are extensive.  Three exact complement-duality
steps reverse the middle transfers and transform the forward final bound
into \(Y_3>1-X\).
\end{proof}

\subsection{T3 endpoint domain}

\begin{proposition}[Symmetry and exact T3 source variables]
\label{prop:s2-t3-source-bridge}
Normalizing the unique center midpoint to \(M_0\), and then choosing the
supported adjacent arm \(r_1\), uses only the \(D_6\)-symmetry of the regular
hexagon.  Those two without-loss-of-generality normalizations are retained.

The later trace-dominating translation is not a \(D_6\)-symmetry.  Its output
is replaced by the exact real parameters
\[
 0<\tau<1,\qquad
 z=\sqrt{1-\tau+\tau^2},\qquad
 0<\beta<\frac{\tau z}{1+z}.
\]
The boundary endpoints and adjacent-radial interval are
\[
 b_-=\beta,\qquad b_+=z-\beta,
\]
\[
 q=1-z+\beta,\qquad
 c_\star=\frac q\tau,\qquad
 u_\star=1+\beta-\tau.
\]
Every support-isolated endpoint parameters has $C_5>1/2$.  If
$C_1\le1/2$, the right endpoint uses the direct linear bound $1-A_1$.  If
$C_1>1/2$, the parameter tuple determines a point of one cell in
Definition~\ref{def:s2-t3-finite-union}.
\end{proposition}

\begin{proof}
The first paragraph is the dihedral normalization in the triangle-forcing and
support-isolation lemmas.  For the non-symmetry step use the exhaustive
translated Type-II equations.  With its original orientation parameter
\(\tau\), one has
\[
 D=\frac1z,\qquad R=\frac{\tau}{z}.
\]
The translated side offset is \(\beta\), and direct substitution gives the
four displayed endpoint formulas.  The center residual orders and the radial
hit/miss orders are precisely the cells
\(\mathcal R_i,\mathcal L_{\varepsilon,j},\mathcal K_h\).

The left endpoint reach lower bound is $C_5=1-a/w>1/2$ by
Proposition~\ref{prop:s2-pure-domain-bounds}.  For the right endpoint, the
branchwise transfer interface gives the direct value $1-A_1$ when
$C_1\le1/2$; this is exactly the existing right-linear row of the complete
$407X$ case analysis.  When $C_1>1/2$, the raw high-radial map supplies one of the
four contact cells, and all defining inequalities of
$\Omega_{\rm T3}$ are satisfied.
\end{proof}

\begin{proposition}[T3 objective bridge]
\label{prop:s2-t3-objective-bridge}
In the high-radial branch $C_1>1/2$, let
$\xi\in\Omega_{\rm T3}$ be the point produced above.  Then
\[
 \mathrm{Cap}(\mathbf e_{\rm T3})-1
 =\Psi_{\rm T3}(\xi).
\]
The complementary branch $C_1\le1/2$ is already the direct right-linear case
and is not part of the optimization domain.
\end{proposition}

\begin{proof}
The variables \(A_1,A_5\) are the far-side uncovered suffix boundary reach lower bounds, while
\(C_1,C_5\) are the complementary radial lower bounds.  On the high-radial branch,
apply \(F(c,x)=g_c(1-x)\) to both endpoint terms.
\end{proof}

\subsection{Adjacent-adjacent exceptional triangle domains}

\begin{proposition}[First adjacent exceptional triangle source bridge]
\label{prop:s2-sc-t3-bridge}
The exact rational variables of the T3-like local profile satisfy
\[
 0\le\theta\le2-\sqrt3,\qquad
 x_L(\theta)\le x\le\frac12-\theta,
\]
and produce
\[
 a=x\rho(\theta),\qquad
 c=x+\theta,\qquad
 u=1-\rho(\theta)+a.
\]
Hence the parameter tuple belongs to \(\Omega_{\rm SC}^{T}\).
\end{proposition}

\begin{proof}
Use
\[
 \theta=\frac{D-1}{\varrho},
 \qquad
 x=\frac{aD}{\varrho}.
\]
The exact inverse formulas for \(D,\varrho\) give the graph identities and
the stated interval bounds.  This replaces the genuine translation reduction;
it does not alter the preceding \(D_6\)-normalization.
\end{proof}

\begin{proposition}[Second adjacent exceptional triangle source bridge]
\label{prop:s2-sc-vd1-bridge}
The Vd1 corner variables satisfy the equations and inequalities of
\(\Omega_{\rm SC}^{V}\).  In particular, absence of a positive interval on
the second adjacent arm is exactly
\[
 \Delta(t)-a-tb-t\le\frac{1-a}{t+1}.
\]
\end{proposition}

\begin{proof}
The two raw endpoints on that arm are
\[
 \lambda_5=\frac{1-a}{t+1},
 \qquad
 \mu_5=\Delta(t)-a-tb-t.
\]
Since \(\lambda_5>0\), absence of a positive interval is equivalent to
\(\mu_5\le\lambda_5\).  The remaining equations are the first-arm graph and
the neighboring-midpoint inequality.
\end{proof}

\subsection{Vd radial domains}

\begin{proposition}[Explicit residual bridge]
\label{prop:s2-vd-residual-bridge}
For a center interval \([L,U]\), the component endpoint and residual used in
the Vd placement proofs are exactly
\[
 \mathsf E(L,U;x),\qquad \mathsf R(L,U;x).
\]
Thus every center-assisted boundary residual belongs to one of the three
explicit order cells in Problem S2-VD.  Moreover, with
\[
 x=\frac{k}{w},\qquad y=\frac{k}{r},
\]
the two center endpoint pairs satisfy
\[
 x^2+x(w+d)+(w+d)^2\le1,\qquad
 y^2+y(r+a)+(r+a)^2\le1.
\]
These are the explicit center-distance constraints included in the
nonadjacent Vd cells.
\end{proposition}

\begin{proof}
If \(x<L\), the initial component ends at \(x\).  If \(L\le x\le U\), it
continues through the interval to \(U\).  If \(U<x\), it already extends to
\(x\).  The cases are disjoint after assigning equality to the middle cell.  The two
quadratic inequalities are the endpoint-distance inequalities for the two
normalized center traces, after substituting the signed-center formulas.
\end{proof}

\begin{proposition}[Explicit corner-graph bridge]
\label{prop:s2-vd-corner-bridge}
Let \(p_1,q_1\) be the exact boundary reaches of an adjacent Vd triangle and let
\(t>0\) be its corner parameter.  The far endpoint of a supported adjacent
interval is
\[
 u_+=\frac{\sqrt{t^2+t+1}-p_1-tq_1-1}{t}.
\]
For a nonadjacent triangle with boundary reaches \(p_v,q_v\), its radial
endpoint is
\[
 C_v=\frac{\sqrt{t^2+t+1}-p_v-tq_v}{t+1}.
\]
Consequently the adjacent and nonadjacent parameter tuples lie in the cells of
Definitions~\ref{def:s2-vd-adjacent-domain} and
\ref{def:s2-vd-nonadjacent-domain}.
\end{proposition}

\begin{proof}
Both identities are exact endpoint formulas of the corner normal form.
Positive support is the explicit order
\(\lambda_+<u_+\); absence of positive support is \(u_+\le\lambda_+\).
The center exits for the three nonadjacent indices are
\[
 d,\qquad \min\left\{\frac ar,\frac dw\right\},\qquad a,
\]
which are exactly the finite order cells in the definition.
\end{proof}

\begin{remark}
The Vd1--supercritical two-chart replacement is not represented by a scalar
problem.  It changes geometric triangles while preserving skeleton coverage and
remains a geometric theorem.
\end{remark}
